% ============================================================
% Stochastic Equations in Classical Physics
% Book chapter. Compiles standalone with pdflatex.
% To fold into a larger manuscript: drop the preamble/\begin{document}
% lines and \input{} this file's body after \chapter{...}.
% ============================================================
\documentclass[11pt]{book}

\usepackage[utf8]{inputenc}
\usepackage[T1]{fontenc}
\usepackage{amsmath,amssymb,amsthm,mathtools,bm}
\usepackage{enumitem}
\usepackage[margin=1in]{geometry}
\usepackage{hyperref}

\newtheorem{remark}{Remark}[chapter]
\theoremstyle{definition}
\newtheorem{example}{Example}[chapter]

\newcommand{\dd}{\mathrm{d}}
\newcommand{\kB}{k_{B}}
\newcommand{\avg}[1]{\left\langle #1 \right\rangle}

\title{Stochastic Equations in Classical Physics}
\author{}
\date{}

\begin{document}

\chapter{Stochastic Equations in Classical Physics}
\label{ch:stochastic}

\section{Determinism, ignorance, and the origin of noise}
\label{sec:intro}

Classical mechanics, in the form developed in the preceding chapters, is
deterministic: given a complete set of initial conditions, Hamilton's
equations (or Newton's law) fix the trajectory of a system for all time.
This determinism is a property of the \emph{ideal world} of the theory --
the mathematical model built from a Lagrangian or a Hamiltonian, a phase
space, and an initial point in it. It is emphatically not a property that
we can exploit in practice for most real systems, because the real world
almost never hands us a closed, low-dimensional set of degrees of freedom.

Consider a grain of pollen suspended in water. A complete Newtonian
description would have to include the grain together with every one of the
$\sim 10^{23}$ water molecules buffeting it. Such a description is exact,
deterministic, and completely useless: we neither know the initial
conditions of the water molecules nor care about their individual
trajectories. What we want is an equation of motion for the grain alone,
and there is no exact, closed equation of that kind -- the grain's dynamics
depends on the full microstate of the fluid.

Stochastic equations are the standard resolution of this impasse. We keep a
small number of "slow" or "relevant" degrees of freedom (here, the position
and velocity of the grain) and replace the influence of everything we have
discarded by an effective random force. The word "random" needs a caveat
that will recur throughout this chapter: the force is not random in the
sense that it violates the underlying determinism of the real world. It is
random in the sense that it is unpredictable \emph{to us}, given the
information we have chosen to retain. A stochastic differential equation is
an ideal-world object -- a specific, well-defined piece of mathematics --
that is deliberately built to model our real-world ignorance of the
eliminated degrees of freedom. This is a somewhat different maneuver from
the ones we met when discussing Hertz's forceless mechanics or nonholonomic
constraints, where the tension was between two competing ideal-world
descriptions of the same real system. Here the tension is built into the
method itself: we manufacture an ideal-world fiction (a random force with
prescribed statistics) precisely because it lets us describe the real world
better than a naive, exact, deterministic account would, given what we
actually know. Whether this maneuver is applied correctly, and what happens
when its idealizations are taken too literally, will be a recurring theme.

\section{Brownian motion: from observation to atoms}
\label{sec:brownian}

The phenomenon that gave the subject its name was reported by the botanist
Robert Brown in 1827: pollen grains suspended in water perform an
incessant, irregular jiggling motion, visible under a microscope, that
never damps out. Brown ruled out a biological origin (the same motion is
seen with inorganic dust), but no dynamical explanation was available for
decades.

The explanation, when it came in 1905 from Einstein (and independently
around the same time from Smoluchowski), was statistical rather than
dynamical in the traditional sense. Einstein did not attempt to write an
equation of motion for the grain. Instead he asked for the probability
$P(x,t)$ that a grain starting at the origin is found near $x$ after time
$t$, given that it receives a huge number of independent, uncorrelated
molecular impacts per unit time. A short argument based on the
Chapman--Kolmogorov idea (the displacement over $[0,t+\tau]$ is the
independent sum of the displacements over $[0,t]$ and $[t,t+\tau]$) shows
that $P$ must obey the diffusion equation
\begin{equation}
\frac{\partial P}{\partial t} = D\,\frac{\partial^2 P}{\partial x^2},
\label{eq:diffusion}
\end{equation}
with diffusion constant $D$ related to the microscopic friction the grain
experiences through the \emph{Stokes--Einstein relation}
\begin{equation}
D = \frac{\kB T}{6\pi \eta a},
\label{eq:stokes-einstein}
\end{equation}
where $a$ is the grain radius, $\eta$ the fluid viscosity, $T$ the
temperature, and $\kB$ Boltzmann's constant. Equation~\eqref{eq:diffusion}
predicts $\avg{x^2(t)} = 2Dt$: the mean-square displacement grows
\emph{linearly} in time, not as $t^2$ as it would for free ballistic
motion. This distinctive signature was confirmed quantitatively by Jean
Perrin's experiments (1908--1909), which used the measured value of $D$ to
extract Avogadro's number and are usually credited with settling, for the
working physics community, the reality of atoms and molecules.

\begin{remark}
It is worth pausing on what kind of result this is. A piece of
ideal-world mathematics -- a linear parabolic PDE describing an idealized
random walk -- was used to answer a real-world existence question ("do
atoms exist?") that had resisted direct attack. The stochastic model here
is not merely a convenient approximation; its quantitative match to
experiment is what turned atomism from a useful bookkeeping device into
an empirically established fact. Few ideal-world constructions in this
book carry that much real-world weight.
\end{remark}

\section{The Langevin equation}
\label{sec:langevin}

Einstein's argument works directly with the probability density and never
mentions a trajectory or a force. In 1908 Langevin gave an equivalent, more
mechanical formulation that has become the standard starting point for
stochastic dynamics. Newton's second law for the grain is written as
\begin{equation}
m\dot v(t) = -\gamma v(t) + \xi(t),
\label{eq:langevin}
\end{equation}
where $-\gamma v$ is the ordinary Stokes viscous drag and $\xi(t)$ is a
\emph{random force} representing the net effect of molecular collisions,
with statistics
\begin{equation}
\avg{\xi(t)} = 0, \qquad \avg{\xi(t)\xi(t')} = 2B\,\delta(t-t').
\label{eq:noise-stats}
\end{equation}
The force is taken to be Gaussian, zero-mean, and $\delta$-correlated in
time ("white noise"): its value at time $t$ is uncorrelated with its value
at any other time $t'\neq t$. This is already an idealization, discussed
further in Section~\ref{sec:white-noise}, since the true molecular force
has a correlation time set by the duration of a collision, many orders of
magnitude shorter than any time scale of interest for the grain -- short
enough that replacing it by zero is an excellent approximation, but not
exactly zero.

Equation~\eqref{eq:langevin} is linear and can be solved by the integrating
factor $e^{\gamma t/m}$:
\begin{equation}
v(t) = v(0)\,e^{-\gamma t/m} + \frac{1}{m}\int_0^t e^{-\gamma (t-s)/m}\,\xi(s)\,\dd s.
\label{eq:v-solution}
\end{equation}
Squaring, averaging over the noise, and using \eqref{eq:noise-stats} gives
\begin{equation}
\avg{v^2(t)} = \avg{v^2(0)}\,e^{-2\gamma t/m}
+ \frac{B}{m\gamma}\Bigl(1 - e^{-2\gamma t/m}\Bigr).
\label{eq:v2-t}
\end{equation}
As $t\to\infty$ the grain must reach thermal equilibrium with the fluid, so
equipartition requires $\avg{v^2}_{\rm eq} = \kB T/m$. Comparing with the
$t\to\infty$ limit of \eqref{eq:v2-t} fixes the noise strength:
\begin{equation}
B = \gamma \kB T
\qquad\Longrightarrow\qquad
\avg{\xi(t)\xi(t')} = 2\gamma \kB T\,\delta(t-t').
\label{eq:fdt-langevin}
\end{equation}
This is our first encounter with the \emph{fluctuation-dissipation
theorem}: the strength of the random force is not a free parameter but is
tied to the friction coefficient $\gamma$ and the temperature $T$. We
return to the general statement in Section~\ref{sec:fdt}.

Integrating \eqref{eq:v-solution} once more and computing the mean-square
displacement reproduces Einstein's diffusive law
$\avg{x^2(t)}\to 2Dt$ for $t \gg m/\gamma$, with
$D = \kB T/\gamma$, consistent with \eqref{eq:stokes-einstein} once
$\gamma = 6\pi\eta a$ (Stokes drag) is inserted. The Langevin and Einstein
pictures are two routes -- one following a (stochastic) trajectory, the
other following a probability density -- to the same physical content;
Section~\ref{sec:fpe} makes the equivalence precise.

\section{White noise as an idealization}
\label{sec:white-noise}

The object $\xi(t)$ deserves scrutiny, because it is more singular than it
looks. A stationary Gaussian process with $\avg{\xi(t)\xi(t')} =
2\gamma\kB T\,\delta(t-t')$ has, formally, infinite variance at each
instant ($\delta(0) = \infty$) and is nowhere continuous, let alone
differentiable. No real physical force behaves this way: any actual
process has some finite correlation time $\tau_c$ (for molecular
collisions, of order $10^{-13}\,\mathrm{s}$ or less), and
$\avg{\xi(t)\xi(t')}$ is some narrow but smooth function of $t-t'$, not a
literal $\delta$-function. White noise is the $\tau_c \to 0$ limit of this
family of processes -- a genuine \emph{singular limit} in the same sense
encountered earlier in this book when relating classical mechanics to its
parent theories: the limiting object (a $\delta$-correlated process) has
qualitatively different properties (non-differentiable sample paths, an
ambiguity discussed in Section~\ref{sec:ito}) from every process in the
family that approaches it.

The right way to think about $\xi(t)$ is not as a function at all but as a
random \emph{distribution} (generalized function): it is defined by the
values of its integrals $\int f(t)\xi(t)\,\dd t$ against smooth test
functions $f$, and only these integrated quantities are guaranteed to be
finite, ordinary random variables. This is formalized through the
\emph{Wiener process} $W(t)$, defined as the Gaussian process with
\begin{equation}
W(0)=0,\qquad \avg{W(t)} = 0, \qquad \avg{\bigl(W(t)-W(s)\bigr)^2} = |t-s|,
\end{equation}
and independent increments. $W(t)$ is continuous but almost surely
nowhere differentiable; $\xi(t)$ is the (nonexistent, but formally useful)
derivative $\dd W/\dd t$. Equation~\eqref{eq:langevin} is more properly
written as the \emph{stochastic differential equation} (SDE)
\begin{equation}
m\,\dd v = -\gamma v\,\dd t + \sqrt{2\gamma \kB T}\;\dd W(t),
\label{eq:sde-langevin}
\end{equation}
an integral equation in disguise: $v(t) - v(0) = \int_0^t(\ldots)\dd s +
\sqrt{2\gamma\kB T}\bigl(W(t)-W(0)\bigr)$, with all the content in the last,
perfectly well-defined stochastic integral.

\section{The It\^o--Stratonovich ambiguity}
\label{sec:ito}

Equation~\eqref{eq:sde-langevin} is \emph{additive}: the noise term
$\sqrt{2\gamma\kB T}\,\dd W$ does not multiply any function of $v$, so there
is no subtlety in defining the corresponding stochastic integral. Trouble
appears as soon as the noise couples to the state variable, as it does in
countless applications (a diffusion coefficient that depends on position,
multiplicative noise in population models, noise entering through a
state-dependent coupling). Consider a general one-dimensional SDE
\begin{equation}
\dd x = a(x,t)\,\dd t + b(x,t)\,\dd W(t).
\label{eq:general-sde}
\end{equation}
Because $W(t)$ has unbounded variation on every interval, the Riemann-sum
definition of $\int b(x,t)\,\dd W$ depends on \emph{where in each
subinterval} $b$ is evaluated -- a dependence that vanishes for ordinary,
finite-variation integrals but survives here. Two conventions dominate the
literature:
\begin{itemize}[leftmargin=1.6em]
\item \textbf{It\^o}: evaluate $b$ at the \emph{left} endpoint of each
subinterval. This makes $b(x(t),t)$ statistically independent of the noise
increment $\dd W(t)$ that follows it, which is convenient for probability
theory (in particular, $\avg{f(x(t))\,\dd W(t)} = 0$ exactly) but means
ordinary calculus (the chain rule) fails; it is replaced by It\^o's lemma,
which carries an extra second-derivative term.
\item \textbf{Stratonovich}: evaluate $b$ at the \emph{midpoint} of each
subinterval. This convention obeys the ordinary chain rule, at the cost of
$b(x(t),t)$ no longer being independent of the subsequent noise increment.
\end{itemize}
The two conventions give \emph{different} solutions $x(t)$ for the same
formal equation \eqref{eq:general-sde} whenever $b$ depends on $x$. They are
related by a well-known conversion formula: the Stratonovich equation
$\dd x = a\,\dd t + b\circ \dd W$ is equivalent to the It\^o equation
\begin{equation}
\dd x = \left(a + \tfrac{1}{2} b\,\frac{\partial b}{\partial x}\right)\dd t + b\,\dd W.
\label{eq:ito-strat}
\end{equation}

This is not a mere technical footnote; it is a physical question in
disguise. Equation~\eqref{eq:general-sde} is, after all, itself an
idealization -- the $\tau_c\to 0$ limit of a system driven by smooth,
rapidly fluctuating \emph{colored} noise, exactly as in
Section~\ref{sec:white-noise}. The Wong--Zakai theorem answers the
question of which convention is "physical": if $\xi(t)$ is approximated by
a family of smooth Gaussian processes with correlation time $\tau_c$, and
one solves the resulting \emph{ordinary} differential equation (well
defined, since the noise is smooth) and then takes $\tau_c\to 0$, the
limit converges to the \textbf{Stratonovich} solution, not the It\^o one.
Physical multiplicative noise -- noise that is really a $\tau_c\to0$ limit
of some smooth real-world process -- should therefore be interpreted in
the Stratonovich sense; the It\^o convention is best reserved for
situations, common in finance and in discrete-time models of chance, where
the "left-endpoint, no look-ahead" property is itself part of the
real-world content being modeled, rather than an artifact of the noise's
idealized whiteness.

\begin{remark}
The It\^o--Stratonovich ambiguity is a clean illustration of a hazard
specific to this chapter's method. Ordinary differential equations, once
written down, have no such interpretational fork: the ideal-world object
and its solution are unambiguous. The moment we idealize a real
fluctuating force as literally $\delta$-correlated, we introduce a fork in
the road that has no counterpart in the pre-limit, physically honest
description. Failing to ask which branch corresponds to the real-world
limit -- treating white noise as simply "there" rather than as a stand-in
for a limit of smooth processes -- is exactly the kind of real-world/
ideal-world conflation this book keeps returning to, here in a
comparatively elementary and very concrete form.
\end{remark}

\section{The Fokker--Planck equation}
\label{sec:fpe}

Einstein's diffusion equation and Langevin's stochastic trajectory are two
sides of the same coin; the general bridge between an SDE and a PDE for
its probability density is the \emph{Fokker--Planck equation} (also called
the forward Kolmogorov equation). For the general one-dimensional It\^o
SDE \eqref{eq:general-sde}, the probability density $P(x,t)$ of $x(t)$
obeys
\begin{equation}
\frac{\partial P}{\partial t}
= -\frac{\partial}{\partial x}\bigl[a(x,t)P\bigr]
+ \frac{1}{2}\frac{\partial^2}{\partial x^2}\bigl[b(x,t)^2 P\bigr].
\label{eq:fpe-general}
\end{equation}
This can be derived either directly from the Chapman--Kolmogorov equation
for the transition probability (a short-time expansion known as the
Kramers--Moyal expansion, truncated at second order because the noise is
Gaussian) or, equivalently, from It\^o's lemma applied to an arbitrary
test function together with an integration by parts. The first term on the
right, the \emph{drift} term, transports probability along the
deterministic part of the flow; the second, the \emph{diffusion} term,
spreads it out.

Two special cases connect back to what we have already derived:
\begin{itemize}[leftmargin=1.6em]
\item Applying \eqref{eq:fpe-general} to the velocity SDE
\eqref{eq:sde-langevin} ($a = -\gamma v/m$, $b^2 = 2\gamma\kB T/m^2$) gives
the \emph{Klein--Kramers equation for the velocity distribution},
\begin{equation}
\frac{\partial P(v,t)}{\partial t}
= \frac{\gamma}{m}\frac{\partial}{\partial v}\bigl[vP\bigr]
+ \frac{\gamma \kB T}{m^2}\frac{\partial^2 P}{\partial v^2},
\label{eq:ou-fpe}
\end{equation}
whose stationary solution is the Maxwell--Boltzmann distribution
$P_{\rm eq}(v)\propto \exp(-mv^2/2\kB T)$, and whose full time-dependent
solution (for a Gaussian initial condition) defines the
\emph{Ornstein--Uhlenbeck process} -- historically the first stochastic
process, after the Wiener process itself, to be solved in closed form
(Uhlenbeck and Ornstein, 1930).
\item In the strongly damped (\emph{overdamped} or \emph{Smoluchowski})
regime, appropriate whenever $\gamma/m$ is much faster than any other rate
in the problem, inertia can be dropped from \eqref{eq:langevin} entirely
($m\dot v \to 0$), leaving $\gamma \dot x = F(x) + \xi(t)$ for a particle in
an external force $F(x) = -U'(x)$. The corresponding Fokker--Planck
equation for the position alone is the \emph{Smoluchowski equation}
\begin{equation}
\frac{\partial P(x,t)}{\partial t}
= \frac{\partial}{\partial x}\!\left[\frac{U'(x)}{\gamma}P\right]
+ D\,\frac{\partial^2 P}{\partial x^2}, \qquad D = \frac{\kB T}{\gamma},
\label{eq:smoluchowski}
\end{equation}
which reduces to Einstein's bare diffusion equation \eqref{eq:diffusion}
when $U\equiv 0$, and is the workhorse equation for Brownian motion in a
potential (Section~\ref{sec:kramers}).
\end{itemize}

\section{The fluctuation--dissipation theorem}
\label{sec:fdt}

The relation $B = \gamma \kB T$ found in \eqref{eq:fdt-langevin} is a
special case of a much more general statement. Whenever a system is
coupled linearly to a large thermal reservoir, the reservoir does two
things at once: it damps the system's motion (dissipation, characterized
by a friction kernel or, in the memoryless case, a single coefficient
$\gamma$) and it drives the system with a fluctuating force (characterized
by a noise spectrum). These are not independent effects with independently
adjustable strengths -- both arise from the same microscopic coupling to
the same reservoir at the same temperature, and the \emph{fluctuation-
dissipation theorem} states that the noise spectrum and the dissipation
kernel are related by a universal formula fixed by $T$ alone. In the
simplest, memoryless (Markovian) case this collapses to
\eqref{eq:fdt-langevin}; more generally, for a friction kernel
$\gamma(t-t')$ (memory friction), the theorem reads
\begin{equation}
\avg{\xi(t)\xi(t')} = \kB T\,\gamma(t-t'),
\end{equation}
with the memoryless Langevin equation recovered when $\gamma(t-t') =
2\gamma\,\delta(t-t')$. The same relation, in its electrical-circuit guise,
is the origin of Johnson--Nyquist noise: a resistor of resistance $R$ at
temperature $T$ develops a fluctuating open-circuit voltage with spectral
density $4\kB T R$, exactly as required for the resistor's Joule
dissipation and its thermal voltage noise to be two faces of one coupling
to the electromagnetic and lattice degrees of freedom it is built from.

\begin{remark}
The fluctuation-dissipation theorem is what keeps the ideal-world
construction of Section~\ref{sec:langevin} honest. Nothing in the bare
form of the Langevin equation \eqref{eq:langevin} forces $B$ and $\gamma$
to be related; one could simply posit \emph{some} friction and \emph{some}
independent noise strength, and the mathematics of Sections
\ref{sec:white-noise}--\ref{sec:fpe} would go through unchanged. Doing so,
however, would in general drive the stationary distribution away from the
Boltzmann distribution $e^{-U/\kB T}$, in violation of the second law: an
uncoupled, freely chosen $(\gamma,B)$ pair generically describes a system
being pumped by its own noise into a non-equilibrium steady state,
extracting or dumping work for nothing. The fluctuation-dissipation
relation is precisely the extra real-world input -- equilibrium
statistical mechanics -- needed to keep the stochastic ideal-world model
from describing a thermodynamically impossible engine. It is a strong
example of a constraint that does not follow from the internal mathematics
of the model at all, and must be imported from outside it.
\end{remark}

\section{Master equations for discrete states}
\label{sec:master}

Not every classical stochastic system is naturally described by a
continuous variable diffusing under noise. Radioactive decay, chemical
reaction kinetics, and population or queueing models are more naturally
posed on a discrete (often countable) set of states $n$, with transitions
occurring as instantaneous jumps at random times. The governing equation
is the \emph{master equation},
\begin{equation}
\frac{\dd P_n(t)}{\dd t} = \sum_{m\neq n}\Bigl[W(n\mid m)\,P_m(t) - W(m\mid n)\,P_n(t)\Bigr],
\label{eq:master}
\end{equation}
where $W(n\mid m)$ is the transition rate from state $m$ to state $n$.
Equation~\eqref{eq:master} is a direct statement of probability
conservation applied to gain and loss terms, and it is Markovian in the
same sense as the Fokker--Planck equation: the future depends on the
present state alone, not on the history that produced it.

The master equation and the Fokker--Planck equation are not competing
descriptions but limits of a single, more general framework (the full
Kramers--Moyal expansion of the Chapman--Kolmogorov equation). When the
individual jumps are small compared to the scale over which $P$ varies --
as happens, for instance, for a chemical system with a very large number
of molecules, where each reaction event changes the population by a
single, comparatively negligible unit -- the master equation's sum can be
Taylor-expanded in the jump size, reproducing a Fokker--Planck equation
with drift and diffusion coefficients built from the first two jump
moments (the van Kampen system-size expansion). Truncating consistently
at second order is what turns a fundamentally discrete, jump-driven
process into an effectively continuous diffusion; higher-order terms,
dropped in this limit, are the trace of the discreteness left behind.

\section{Kramers' escape problem}
\label{sec:kramers}

A classic application that draws on nearly everything above is the rate at
which thermal noise drives a particle out of a metastable potential well.
Let $U(x)$ have a local minimum at $x_a$ and a barrier at $x_b > x_a$ of
height $\Delta U = U(x_b) - U(x_a) \gg \kB T$, and let the particle obey
the overdamped Langevin equation of Section~\ref{sec:fpe}. In the absence
of noise the particle sits at $x_a$ forever; with noise, it occasionally
receives a large enough kick to cross the barrier, and the mean time
$\tau$ spent in the well before escaping can be computed from the
stationary Smoluchowski equation \eqref{eq:smoluchowski} with an absorbing
boundary placed just past $x_b$ (a standard \emph{mean first-passage time}
calculation). The result, worked out by Kramers in 1940, is an
Arrhenius-type law
\begin{equation}
\frac{1}{\tau} \;\approx\; \frac{\sqrt{U''(x_a)\,|U''(x_b)|}}{2\pi\gamma}\;
e^{-\Delta U/\kB T},
\label{eq:kramers-rate}
\end{equation}
valid in the high-friction regime, with a friction-independent prefactor
correction in the low-friction (energy-diffusion-limited) regime that
Kramers also worked out. The exponential factor $e^{-\Delta U/\kB T}$ is
the familiar Boltzmann-weighted Arrhenius law of chemical kinetics; what
the stochastic calculation supplies, beyond phenomenological chemistry, is
the prefactor and its explicit dependence on the local curvature of the
potential and on the friction $\gamma$ -- information invisible to
equilibrium thermodynamics alone, since escape is fundamentally a dynamical,
not merely a thermodynamic, question. Kramers' problem remains the
standard reference point for noise-activated transitions across
chemical reaction rate theory, nucleation, and the switching of bistable
systems generally.

\section{Closing remarks: reading the dictionary both ways}
\label{sec:closing}

It is worth collecting, in one place, the real-world/ideal-world
correspondences this chapter has relied on, since they run in a direction
slightly different from the ones met earlier in this book.

\begin{itemize}[leftmargin=1.6em]
\item In the chapters on classical mechanics proper, paradoxes typically
arose because an ideal-world construction (a coordinate system, a
constraint, a limiting procedure) was mistaken for -- or failed to
correspond to -- some feature of the real system it was meant to describe.
The direction of the error was: too much confidence that the mathematics
faithfully tracked the physics.
\item Here the situation is reversed. We \emph{start} from the acknowledged
fact that we cannot track the real system in full (too many degrees of
freedom, no access to their initial conditions) and \emph{deliberately}
introduce an ideal-world fiction -- a $\delta$-correlated random force,
a Wiener process, white noise -- that has no direct real-world referent at
all. The risk is not overconfidence in the mathematics but the opposite:
forgetting that the fiction is a fiction, and that it is licensed only as
a limit of an honest, non-singular, real-world process (colored noise
with a small but nonzero correlation time). Section~\ref{sec:ito} showed
concretely what goes wrong when this is forgotten: the naive ideal-world
object (white noise multiplying a state-dependent coefficient) admits two
inequivalent completions, and only one of them, the Stratonovich one, is
the actual limit of the real-world family that white noise is supposed to
stand in for.
\item The fluctuation-dissipation theorem (Section~\ref{sec:fdt}) is the
clearest single illustration of the discipline this requires: the
ideal-world model (friction coefficient here, noise amplitude there) has
free parameters that look independent on the page but are not free at all
once the model is required to reproduce a real-world constraint
(thermal equilibrium, the second law) that has no analogue inside the
stochastic differential equation itself.
\end{itemize}

Stochastic equations, in short, are not a retreat from the deterministic,
mechanistic picture of the earlier chapters but its natural extension: a
disciplined way of being honest about what we do not track, provided we
remember, at each of the junctures above, that the discipline is imported
from the real world and does not follow automatically from the ideal-world
mathematics alone.

\end{document}
